% !TEX root = userguide.tex

\def\headdate{11th May 2025}

\section{Viscoelastic fluid flow using implicit time discretization schemes}
\label{sec:implicit}

In this section some examples will be given that solve the equations 
for viscoelastic fluid flow using implicit time discretization schemes.

We start with the fully coupled case where
all variables (velocity, pressure, conformation) are implicit.
Then we show the stress-explicit and stress-implicit decoupled schemes, where
the constitutive equation is discretized in time using fully implicit schemes
with respect to the conformation variables only.

Advantages of implicit time discretization schemes in \tfem, as compared to explicit ones, include
\begin{itemize}
  \item Larger time steps are usually permitted before numerical instabilities set in.
  \item The Jacobian matrix is available and can be used to perform
      an eigenvalue analysis and obtain linear stability information.
  \item Both the standard conformation formulation and the CDT formulation (see Sec.~\ref{sec:CDTformulation}) are available. 
\end{itemize}
Disadvantages and/or limitations of implicit time discretization schemes in \tfem\ are:
\begin{itemize}
  \item Each time step is more expensive, because we need iteration (Newton-Raphson) to solve the
nonlinear equations and the matrices to solve are (sometimes much) larger.
  \item The CDT formulation (see Sec.~\ref{sec:CDTformulation}) is inherently time dependent, except fur pure extensional flows. This means, that both time and space discretization errors are present. Even if you are mainly interested in a steady state solution, increasing the time step will also increase the error in the steady state.
  \item Implicit time discretization for the LCR (log-conformation representation) is not yet available in \tfem.
\end{itemize}

\subsection{Implicit time discretization schemes}
\label{sec:implicitschemes}

Let's recall the equations for the flow of a single-mode incompressible viscoelastic fluid
in a region $\Omega$ (see Eqs.~\eqref{eq:ve1}--\eqref{eq:ve3}):
\begin{alignat}{2}
  \rho(\pderiv{\vek u}t+\vek u\cdot\nabla\vek u)
      -\nabla\cdot(2\eta_\text{s}\ten D(\vek u))+\nabla p-\nabla\cdot\ten \tau(\ten c)
        &=\rho\vek b,\qquad&&\text{in }\Omega\label{eq:impve1}\\
\nabla\cdot\vek u&=0,\qquad&&\text{in }\Omega\\
        \pderiv{\ten c}{t}
+\vek u\cdot\nabla\ten c-\ten L\cdot\ten c-\ten c\cdot\ten L^T
 +\ten f(\ten c)&=\ten 0,\qquad&&\text{in }\Omega\label{eq:impve3}
\end{alignat}
 In these equations $\vek u$ is the velocity, $p$ is the
pressure and $\ten c$ is the conformation tensor, which are the unknowns to be
solved as a function of position $\vek x$ and time $t$. Furthermore $\rho$ is
the density of the fluid, $\eta_\text{s}$ is the solvent viscosity and
 $\ten D=(\nabla\vek u+\nabla\vek u^T)/2$, with $\ten L=(\nabla\vek u)^T$ and
$\vek b$ is a body force per unit mass (e.g.\ gravity).
The polymer stress tensor $\ten\tau(\ten c)$ and the tensor function 
$\ten f(\ten c)$ are 
functions of the conformation tensor and depend on the model (see
Appendix~\ref{chap:viscoelasticmodels}).

We ignore the inertia terms for now. Applying an implicit Euler (BDF1) time discretization to Eq.~\eqref{eq:impve3} and evaluating the momentum balance and continuity equation at the new time $t_{n+1}$, we obtain the \emph{fully coupled scheme}:
\begin{alignat*}{2}
        -\nabla\cdot(2\eta_\text{s}\ten D(\vek u_{n+1}))+\nabla p_{n+1}-\nabla\cdot\ten \tau(\ten c_{n+1})
        &=\rho\vek b,\qquad&&\text{in }\Omega\\
\nabla\cdot\vek u_{n+1}&=0,\qquad&&\text{in }\Omega\\
        \frac{\ten c_{n+1}-\ten c_n}{\Delta t}
+\vek u_{n+1}\cdot\nabla\ten c_{n+1}-\ten L(\vek u_{n+1})\cdot\ten c_{n+1}-\ten c_{n+1}\cdot\ten L^T(\vek u_{n+1})
 +\ten f(\ten c_{n+1})&=\ten 0,\qquad&&\text{in }\Omega\label{eq:dimpve3}
\end{alignat*}
where $\Delta t=t_{n+1}-t_n$.
\newpage
Remarks:
\begin{enumerate}
\item The system is nonlinear and iteration, like Newton-Raphson, is needed to solve it.
\item It is a fully coupled system between the variables $(\vek u_{n+1},p_{n+1},\ten c_{n+1})$ with a Jacobian of the form
\[
   \begin{pmatrix}
     S & L & B_1 \\
     L^T & 0 & 0 \\
     C_1 & 0 & D_1
    \end{pmatrix}
\]
which extends to
\[
   \begin{pmatrix}
     S & L & B_1 & B_2 & \cdots & B_m \\
     L^T & 0 & 0 & 0 & \cdots & 0 \\
     C_1 & 0 & D_1 & 0 & \cdots & 0 \\
     C_2 & 0 & 0 & D_2 & \cdots & 0 \\
      \vdots  & \vdots & \vdots & \vdots & \ddots & \vdots \\
     C_ m& 0 & 0 & 0& \cdots & D_m \\
    \end{pmatrix}
\]
assuming $m$ modes are used in the constitutive equation.
The matrix is sparse, even on element level, and \tfem\ can take that into account to avoid unnecessary filling of the system matrix (see Sec.~\ref{sec:excludezeroblocks}).
\item If the viscoelastic modes are coupled, like in the EGP model with a global von-Mises stress (see Sec.~\ref{sec:egp_flow}), the `\textit{D}' matrix becomes a full matrix:
\begin{equation}
    \begin{pmatrix}
     S & L & B_1 & B_2 & \cdots & B_m \\
     L^T & 0 & 0 & 0 & \cdots & 0 \\
     C_1 & 0 & D_{11} & D_{12} & \cdots & D_{1m} \\
     C_2 & 0 & D_{21} & D_{22} & \cdots & D_{2m} \\
      \vdots  & \vdots & \vdots & \vdots & \ddots & \vdots \\
     C_ m& 0 & D_{m1} & D_{m2}& \cdots & D_{mm} \\
    \end{pmatrix}\label{eq:fullDmatrix}
\end{equation}
\item When applying the DEVSS-G scheme, an additional variable $\ten G$ is introduced and in the CE the $\ten L$ is replaced with $\ten G$. The Jacobian matrix with variables $(\ten G_{n+1},\vek u_{n+1},p_{n+1},\ten c_{n+1})$ now becomes of the form
\[
   \begin{pmatrix}
     E & F & 0 & 0 \\
    F^T & S & L & B_1 \\
    0 & L^T & 0 & 0 \\
    G_1 & C_1 & 0 & D_1
    \end{pmatrix}
\]
which extends to
\[
   \begin{pmatrix}
     E & F & 0 & 0 & 0 & \cdots & 0 \\
     F^T & S & L & B_1 & B_2 & \cdots & B_m \\
     0 & L^T & 0 & 0 & 0 & \cdots & 0 \\
     G_1 & C_1 & 0 & D_1 & 0 & \cdots & 0 \\
     G_2 & C_2 & 0 & 0 & D_2 & \cdots & 0 \\
      \vdots &\vdots  & \vdots & \vdots & \vdots & \ddots & \vdots \\
     G_m & C_m& 0 & 0 & 0& \cdots & D_m \\
    \end{pmatrix}
\]
assuming $m$ modes are used in the constitutive equation and that these are not coupled.
\item For the fully coupled scheme no special techniques are needed for zero solvent fluids ($\eta_\text{s}=0$). Hence the stress-implicit scheme of Sec.~\ref{sec:implicitstress} is not applicable here.
\item Other time discretization methods (Crank-Nicolson, BDF2, $\theta$-method)
 can be easily applied.
 \item Since all unknowns are solved `in one go', the main program is shorter and easier to read than decoupled methods.
\item The Jacobian can be used to perform a linear stability analysis. See Sec.~\ref{sec:linearstability}.
\item The Jacobian can be used to perform a Newton-Raphson iteration to find a steady state solution. See Sec.~\ref{sec:steadyviscoelastic}.
\end{enumerate}

The fully coupled scheme can be very expensive, in particular with multiple mode models. Therefore, a decoupled scheme can be more attractive in that case.
Again applying an implicit Euler (BDF1) time discretization to Eq.~\eqref{eq:impve3}, evaluating the momentum balance and continuity equation at the new time $t_{n+1}$, however introducing a (known) velocity $\hat{\vek u}_{n+1}$,
we get the \emph{decoupled scheme}:
\begin{alignat*}{2}
        -\nabla\cdot(2\eta_\text{s}\ten D(\vek u_{n+1}))+\nabla p_{n+1}-\nabla\cdot\ten \tau(\ten c_{n+1})
        &=\rho\vek b,\qquad&&\text{in }\Omega\\
\nabla\cdot\vek u_{n+1}&=0,\qquad&&\text{in }\Omega\\
        \frac{\ten c_{n+1}-\ten c_n}{\Delta t}
+\hat{\vek u}_{n+1}\cdot\nabla\ten c_{n+1}-\ten L(\hat{\vek u}_{n+1})\cdot\ten c_{n+1}-\ten c_{n+1}\cdot\ten L^T(\hat{\vek u}_{n+1})
 +\ten f(\ten c_{n+1})&=\ten 0\qquad&&\text{in }\Omega\label{eq:ddimpve3}
\end{alignat*}
Remarks:
\begin{enumerate}
\item If we apply the \emph{explicit stress} formulation (see Sec.~\ref{sec:explicitstress}), the velocity vector 
$\hat{\vek u}_{n+1}$ is obtained from prediction using previous time steps. The conformation $\ten c_{n+1}$ can be 
obtained by solving the constitutive equation for known velocity vector 
$\hat{\vek u}_{n+1}$.
The momentum balance and continuity equation are subsequently solved for obtaining $(\vek u_{n+1},p_{n+1})$ by substitution of $\ten c_{n+1}$. This only works for $\eta_\text{s}\neq0$.
\item If we apply the \emph{implicit stress} formulation (see Sec~\ref{sec:implicitstress}), we first solve for 
$(\vek u_{n+1},p_{n+1})$ and obtain the conformation $\ten c_{n+1}$ by substituting $\hat{\vek u}_{n+1}=\vek u_{n+1}$ in the constitutive equation. This also works for $\eta_\text{s}=0$.
\item The constitutive equation is a nonlinear equation in terms of $\ten c_{n+1}$ and needs to be solved by iteration, for which we apply the Newton-Raphson method.
\item For multi-mode models, each mode is solved separately (decoupled).
\item Other time discretization methods (Crank-Nicolson, BDF2, $\theta$-method) are available.
\end{enumerate}

\subsection{Planar Couette flow for a Weissenberg number (Wi) of 10}
\label{sec:couette_implicit}

Various Couette flow examples using implicit time integration are available. The constitutive equation is always solved with a fully implicit scheme. However the coupling with the momentum balance is either stress explicit, stress implicit of fully implicit. Furthermore the formulation is either the standard c formulation or the b formulation (logc is not available). The following table shows the examples:\medskip\\
\begin{tabular}{llll}
Example & c or b  & coupling & postprocessing \\\hline
 \texttt{couette17.f90} & c & explicit stress & \texttt{couette\_post17.f90} \\
 \texttt{couette18.f90} & c & implicit stress & \texttt{couette\_post17.f90} \\
 \texttt{couette19.f90} & b & explicit stress & \texttt{couette\_post17.f90} \\
 \texttt{couette20.f90} & b & implicit stress & \texttt{couette\_post17.f90} \\
 \texttt{couette21.f90} & c & fully implicit & \texttt{couette\_post21.f90} \\
 \texttt{couette22.f90} & b & fully implicit & \texttt{couette\_post21.f90} 
\end{tabular}\medskip\\
We refer to Sec.~\ref{sec:couette} for a description of the Couette problem.

Examples \texttt{couette19a.f90}, \texttt{couette20a.f90} and \texttt{couette22a.f90} are similar to examples \texttt{couette19.f90}, \texttt{couette20.f90} and \texttt{couette22.f90}, respectively, but use element level rotation reinitialization instead of the nodal based one.

The examples can be obtained by typing at the
command line:
\medskip
\begin{quote}
   \texttt{getexample couette}
\end{quote}

\subsection{Flow around a cylinder between two plates (implicit; decoupled and fully implicit)}
\label{sec:cylinderimplicit}
The examples can be obtained by typing at the command line:
\begin{quote}
   \texttt{getexample confined\_cylinder\_b}
\end{quote}

The examples describe the flow
around a fixed cylinder positioned between two walls. For the examples \texttt{cylinder3.f90} and \texttt{cylinder4.f90} the numerical scheme is decoupled with
stress-explicit (see Sec.~\ref{sec:explicitstress}) and stress-implicit (see
Sec.~\ref{sec:implicitstress}) approaches for the two examples, respectively. The example \texttt{cylinder5.f90}
shows the use of a fully implicit scheme. The example \texttt{cylinder6.f90} is similar to \texttt{cylinder5.f90} but has the option to reduce the memory footprint of the system matrix when dealing with a large number of modes.

Examples \texttt{cylinder3a.f90}, \texttt{cylinder4a.f90}, \texttt{cylinder5a.f90} and \texttt{cylinder6a.f90} are similar to examples \texttt{cylinder3.f90}, \texttt{cylinder4.f90}, \texttt{cylinder5.f90} and \texttt{cylinder6.f90}, respectively, but use element level rotation reinitialization instead of the nodal based one.
Example \texttt{cylinder9.F90} is similar to \texttt{cylinder6a.f90}, but reads from an input file:
 \begin{quote}
   \texttt{cylinder9 < input\_cylinder9.txt}
\end{quote}
and there is also the option to restart from a previously saved solution. The example \texttt{cylinder9.F90} is used to create the results in \cite{JaenssonHulsen2024}.

A default coarse mesh (\texttt{confined\_cylinder1.msh}) is supplied, but in the
subdirectory \texttt{meshes} files are available to generate more refined 
meshes.

Examples \texttt{cylinder\_c3.f90}, \texttt{cylinder\_c4.f90}, \texttt{cylinder\_c5.f90}, \texttt{cylinder\_c6.f90} and \texttt{cylinder\_c9.f90} are also available
and are similar to \texttt{cylinder3.f90}, \texttt{cylinder4.f90}, \texttt{cylinder5.f90}, \texttt{cylinder6.f90} and \texttt{cylinder9.f90}, respectively. However 
the problems are solved using the conformation tensor approach. Since the log-conformation approach cannot yet be applied to the implicit time discretization, the Weissenberg number that can be imposed is limited for the conformation approach.

